% Рисунок с доски: релаксация напряжения в материале Максвелла при
% e = e_0 H(t); функция релаксации R(t) = 2μ e^{-t/τ}.
\begin{figure}[htbp]
\centering
\begin{tikzpicture}[>={Latex[length=4pt]}]
  \draw[->] (0,0) -- (6.3,0) node[anchor=north east] {$t$};
  \draw[->] (0,0) -- (0,3.1) node[anchor=south east] {$\Sdev$};

  \draw[very thick, blue!60!black, domain=0:6, samples=120, smooth]
    plot (\x, {2.6*exp(-\x/1.5)});
  \node[anchor=east, font=\footnotesize] at (-0.05,2.6) {$2\mu\,\dev_0$};

  \draw[gray] (0,2.6) -- (1.5,0);
  \draw[dashed, gray] (1.5,0) -- (1.5,{2.6*exp(-1)});
  \fill (1.5,{2.6*exp(-1)}) circle (1.2pt);
  \node[anchor=north, font=\footnotesize] at (1.5,-0.05) {$\tau$};
  \node[font=\footnotesize, anchor=west] at (2.6,1.4) {релаксация};
  \draw[->, gray] (2.5,1.3) -- (1.9,0.75);
\end{tikzpicture}
\caption{Релаксация: при мгновенно заданной и далее постоянной деформации
напряжение падает по закону $R(t) = 2\mu e^{-t/\tau}$, $\tau = \eta/\mu$ —
время релаксации.}
\label{fig:relaxation-function}
\end{figure}
